\section{The Second Published Witness: Peter Pernin}

\subsection{Where the account sits}

The second witness is better known and worse understood. Peter Pernin
was the missionary priest at Peshtigo and Marinette who lived through
the firestorm of the night of 8 October 1871 by standing in the
Peshtigo River, and whose memoir of it remains the classic eyewitness
account of that catastrophe. He published it at Montreal in 1874, in
French and in English in the same year, and the Champion material is
not part of the memoir at all: it is an appendix, added after the
narrative ends.\endnote{Peter Pernin, \work{The Finger of God Is There!
or, Thrilling Episode of a Strange Event Related by an Eye-Witness}
(Montreal: printed by John Lovell, St.~Nicholas Street, 1874), title
page and Appendix, printed pp.~94--102. The French original of the same
year is \work{Le doigt de Dieu est l\`a! ou \'Episode \'emouvant d'un
\'ev\`enement \'etrange racont\'e par un t\'emoin oculaire} (Montreal:
Eus\`ebe S\'en\'ecal, imprimeur, 1874), Appendice, printed pp.~91--101.
Both were read 2026-07-25 in Canadian Institute for Historical
Microreproductions scans on the Internet Archive, identifiers
\texttt{cihm\_04805} (English) and \texttt{cihm\_23982} (French). Both
title pages carry the same two statements: publication ``with the
approbation of His Lordship the Bishop of Montreal'' and publication
``for the Church of Our Lady of Lourdes, in Marinette, State of
Wisconsin.'' The English volume prints the approbation in full: the
bishop has read the work, has ``been deeply touched by it,'' and
recommends that the faithful of his diocese keep a copy at home ``so as
to read and re-read it frequently in the family circle.'' The
approbation commends the book; it makes no judgment about the
Robinsonville events.}

Two features of that provenance bear on the appendix's evidential
value, and they pull in opposite directions. The book was a
fundraiser, sold for the benefit of a church at Marinette dedicated to
Our Lady of Lourdes, and its governing theme is the finger of God and
the intervention of the Blessed Virgin; the appendix exists to supply a
second instance of that theme. But Pernin also names his own restraint:
``I must necessarily abridge details and be guarded in my words so as
to avoid wounding the modesty of the personages mentioned in this
recital, many of whom are still living and may later read these
pages.'' A witness who tells you he is compressing has told you how to
read his silences.

He also had access. He was a priest working across Green Bay from the
Belgian settlement; he had spoken with people who knew Adele as a child
in Belgium; and he invites his readers to go and question her
themselves, ``this humble place of pilgrimage, as yet in its infancy,
and the only one, I believe, of the nature in the United
States.''\endnote{Pernin, \work{The Finger of God Is There!}, pp.~95
and 101--102; French, pp.~92 and 100--101 (``qui est le seul encore,
je crois, aux \'Etats-Unis'').}

\subsection{The passage}

\begin{witnesstext}{Peter Pernin, \work{The Finger of God Is There!},
Appendix, printed pp.~95--97 (Montreal, 1874), with the French of the
same year collated at every clause}
Among them dwells an unmarried female of about forty years of age, poor
in point of fortune and physical attractions but rich in grace and
virtue. Her name is Ad\`ele Brisse. \dots{} Ten or twelve years ago,
the colony possessed no resident priest, and this fervent Catholic set
out early every Sunday morning on foot, and walked to the neighboring
parish, seven miles and a-half distant, approached the sacraments,
heard Mass and returned in all haste to resume her duties in her humble
home.\par
She was returning one morning, after having approached the Blessed
Eucharist. All at once she perceived above the middle of the little
path she was following through the woods, a Lady of great beauty and
majestic mien, who stood as if suspended between two trees, bordering
the path. Surprised and greatly moved, though not terrified, she fell
on her knees, uttered a short prayer, then rose to her feet again. She
made no mention to her companions of what she had seen, but they had
witnessed her emotion, divined the cause, and the mysterious
manifestation began to be talked about. The week following, Ad\`ele
Brisse walked as usual to the neighboring parish to fulfil her
religious duties, returning from mass accompanied by a number of her
companions when the same apparition appeared to her at the same place.
\dots{} Still many resolved they would accompany her on her next
journey and judge for themselves, which they accordingly did the
following Sunday. After having confessed, communicated and heard mass,
Ad\`ele Brisse took her homeward way, accompanied this time by a
numerous escort, among whom were several men. \dots{} Arrived at the
place where the Lady had already appeared twice, she showed herself,
for the third time, more majestic and loveable than ever. The girl fell
on her knees, no feeling of fear agitating her breast, but instead, a
sentiment of perfect confidence, and then there ensued between herself
and this majestic Lady, a long conversation of which I will repeat only
what is necessary to the elucidation of my subject. Ad\`ele began:\par
``My good mother, what do you want of me?''\par
``That you should instruct my children,'' was the Lady's reply. ``You
have just received my Son within your breast, and you have done well,
but these poor children receive Him without knowing what they do, and
are growing up in ignorance of their religion. I wish you to instruct
and above all prepare them for their first communion.''\par
``How can I do that, my good mother, I am but a poor ignorant creature
myself?''\par
``Go, and fear nothing; I will help you.''\endnote{Pernin, \work{The
Finger of God Is There!}, pp.~95--97; French, pp.~92--95. Public
domain: printed 1874, author died 1909. Printed pp.~95 and 98--100 were
read at 200~dpi from the scan's page images and the quotations here were
collated against them; the seer's name is set with a grave accent in the
1874 printing (``Ad\`ele Brisse''), and the acute of the
optical-character-recognition text is a scanner artifact. Elsewhere the
recognition text has been corrected only where the scanner plainly
misread a letter within a word; no wording has been supplied,
normalized, or harmonized with the French. The French reads at the
corresponding place:
``\emph{Bonne m\`ere, que voul\'ez-vous de moi.} --- \emph{Je veux, lui
r\'epondit la Dame, que tu instruises mes enfants. Tu viens de recevoir
mon fils, et tu as bien fait, mais ces pauvres enfants le re\c{c}oivent
sans savoir ce qu'ils font, et grandissent dans l'ignorance de la
Religion. Je veux que tu les instruises et surtout, que tu les prépares
\`a leur premi\`ere communion.} --- \emph{Comment ferai-je cela, bonne
m\`ere, je ne suis qu'une pauvre ignorante moi-m\^eme?} ---
\emph{Va, et ne crains rien; je t'aiderai.}'' The two texts agree
substantively throughout the appendix; the divergences located are
noted in the body above and in Appendix~A.}
\end{witnesstext}

Pernin then gives the sequel that Starr also gives, and gives it more
fully: Adele going from village to village through rain, snow and heat;
neither fatigue nor ridicule stopping her; a priest at last found for
the colony, who advises her to raise money for a schoolhouse rather
than wear herself out following the children through the woods; and, by
1874, a school ``capable of containing more than one hundred children
who yearly are prepared therein for their first communion'' and ``a
little chapel built in the place where the Blessed Virgin appeared to
her, and in which is kept as a precious relic, the tree sanctified by
her apparent touch.'' The two wooden buildings stand in an enclosure of
``about six acres of land, a gift bestowed on her, which she in turn
made over to the Bishop of Green Bay,'' fenced with boards, with a
narrow path winding round it along which a solemn procession passes
twice a year ``at fixed epochs'' --- a ``procession which attracts more
than four thousand pilgrims to the neighborhood.'' Five or six young
women had by then joined her.\endnote{Pernin, \work{The Finger of God Is
There!}, pp.~97--99; French, pp.~95--97; printed pp.~98--99 read at
200~dpi. The French has ``l'arbre sur lequel s'est faite l'apparition''
where the English has ``the tree sanctified by her apparent touch,'' and
``six arpents'' where the English has ``about six acres.'' The English
printing has a compositor's error at this point, ``twice very year''
for ``twice every year''; the French reads ``\emph{chaque ann\'ee, \`a
deux \'epoques fixes}.'' Pernin does not name the two dates.}

\subsection{What this witness does and does not carry}

\textbf{The date is a problem, and it is Pernin's own.} ``Ten or twelve
years ago'' from an 1874 imprint gives 1862--1864, and the French is no
vaguer: \term{Il y a 10 ou 12 ans}. Pernin attaches the phrase to the
colony's want of a resident priest and then narrates the apparitions as
following, so his implied date for the events is three to five years
later than the October 1859 that the diocese later fixed. This study
does not attempt to reconcile the two. It records that the second
earliest published witness, writing fifteen years after the received
date and with local access, placed the events in the early 1860s, and
that no act of any competent authority has ever adjudicated the point.

\textbf{The sequence differs from Starr's.} Pernin has three
encounters, all on the return from Mass, on three successive Sundays,
with the crowd growing each week. Starr has the first on a return, the
second on an outward journey some weeks later, and the third the same
afternoon. They cannot both be right, and neither of them is the
sequence the shrine narrates today.

\textbf{The distance differs from Starr's.} Pernin's French gives two
and a half leagues, which his English translates as seven and a half
miles; Starr has eleven. Nothing turns on the discrepancy except the
reminder that neither witness was working from a record.

\textbf{The figure is barely described.} ``A Lady of great beauty and
majestic mien, who stood as if suspended between two trees, bordering
the path.'' The two trees are here, in the earliest witness that
mentions them; their species is not. There is no veil, no crown, no
sash, no hair, no light.

\textbf{The identification is again the seer's, not the visitor's.}
Adele addresses her as ``my good mother'' from the first exchange.
There is no self-naming.

\textbf{The words are Eucharistic and catechetical, and one received
clause appears.} Pernin's version of the commission turns on First
Communion and on children receiving the Lord ``without knowing what
they do,'' which is not in Starr; and it closes with ``Go, and fear
nothing; I will help you,'' which \emph{is} the closing clause of the
message as the shrine and the 2010 decree now give it. That is the one
element of the fuller received wording that this study can trace to a
witness of the 1870s.

\textbf{The silence is declared.} Pernin says he will repeat ``only
what is necessary to the elucidation of my subject.'' His omissions
therefore prove much less than Starr's. If a self-identification as
Queen of Heaven, a command of general confession, and a warning of
punishment were current at Robinsonville in 1874, Pernin had a stated
reason to pass over them --- though it is worth noticing that a warning
of punishment would have suited the theology of his own book, which
reads the Peshtigo fire as a summons to repentance and calls the town
``a modern Sodom in the sense that it can serve as an example to
all.''\endnote{Pernin, \work{The Finger of God Is There!}, p.~93;
French, p.~90. The passage closes his memoir proper, immediately before
the Appendix, with Psalm~2 (Vulgate numbering as he cites it): ``And
now, O ye kings, understand: receive instruction you that judge the
earth.''}

\begin{judgmentbox}{Project synthesis --- what Pernin's declared
abridgement is worth}
This study judges that Pernin's appendix establishes four things
securely and one thing not at all. Securely: that by 1874 the site had
a chapel on the reported spot with the tree preserved in it, a school,
a six-acre enclosure conveyed to the Bishop of Green Bay, two annual
processions drawing thousands, and a small community of women; that the
reported commission was already understood as catechesis directed to
First Communion; that ``Go, and fear nothing; I will help you'' was
already part of the received wording; and that the events were, in
1874, unjudged by ecclesiastical authority, which Pernin says in terms.
Not at all: the date. Pernin's own ``ten or twelve years'' is
incompatible with October 1859, and this study can find no way to
resolve the conflict from the sources it reached. The strongest reply
is that Pernin was writing from conversation, not calculation, and that
a priest recalling a lay woman's biography fifteen years on may
casually shave a few years; that reply is plausible, and it is also
unverifiable, which is why the conflict is recorded here rather than
smoothed away.
\end{judgmentbox}
